\newcommand{\V}{\mathcal{V}}

Similar to toric varieties, let $N := \Hom(M, \mathbb{Z})$, $N_\mathbb{Q} := \Hom(M, \mathbb{Q})$, and we denote the perfect pairing by $\<\cdot, \cdot\> : M_\mathbb{Q} \times N_\mathbb{Q} \to \mathbb{Q}$.

Given a valuation $v : \mathbb{C}(X)^\times \to \mathbb{Q}$, we can define a map $\rho(v)(\lambda) := v(f_\lambda)$ for every character $\lambda \in M$; we have $\rho(v) \in N_\mathbb{Q}$, i.e.,
\begin{equation}
\begin{aligned}
	\rho: \{\text{valuations of }\mathbb{C}(X)\} \to N_\mathbb{Q}
\end{aligned}
\end{equation}
is well-defined.

% valuation cone
\begin{lemdefn}[c.f., e.g., {\cite[\S4]{Brion_var_sph}}]\label[lemdefn]{def:valcone}
	After restriction,
	\begin{equation*}
		\rho : \{G\text{-invariant valuations of }\mathbb{C}(X)\} \to N_\mathbb{Q}
	\end{equation*}
	is injective. Moreover, its image is a co-simplicial convex cone, called \emph{valuation cone} of $X$, denoted by $\V$.
\end{lemdefn}
\begin{remark*}
	In other words, each $G$-invariant valuation of $\mathbb{C}(X)$ is uniquely determined by its values on $B$-semi-invariant rational functions, or equivalently, principal $B$-invariant divisors.
\end{remark*}

As a result, we identify a $G$-invariant valuation $v$ with a vector $\rho(v) \in \V$.

As the cone $\V$ is polyhedral, the following lemma is intuitive.
\begin{lemma}\label[lemma]{lem:sph:DivValEQVal}
	Assume $X$ is a projective spherical $G$-variety, we have $\Val^G_X = \DVal^G_X$.
\end{lemma}
\begin{remark*}
	Recall that, in this paper we assumed valuations are $\mathbb{Q}$-valued. Note that there always exists $G$-invariant $\mathbb{R}$-valuations which are not of the form $c \: \ord_E$ for some prime divisor $E$ over $X$ and $c \in \mathbb{R}_{\geq 0}$; for instance, quasi-monomial ($\mathbb{R}$-)valuations (see, for example, \cite[\S3]{JM_val_asym_inva_seq_i}).
\end{remark*}
\begin{sketchproof}
	According to \cite[Proposition 2, \S2.4]{Brion_var_sph}, there is a toroidal $G$-equivariant birational morphism $\pi: X' \to X$. For every $G$-invariant valuation $v \in \Val^G_X$, there is a subdivision of $\V$ where $\mathbb{Q}_{\geq 0} v \subseteq \V$ is a cone. The subdivision of $\V$ corresponds to birational model $X''$ of $X'$ (thus a birational model of $X$), and the ray $\mathbb{Q}_{\geq 0} v$ defines a prime divisor of $X''$.
\end{sketchproof}

This lemma is stronger than \cite[Proposition 2.19]{Golota_Delta_Fano_large_aut}, which proves the subset $\DVal^G_X$ is dense in the space of real-valued $G$-invariant valuations.


